% !TeX root = ../papers/nc-p3.tex
\section{可逆二平面分解の点スキームと孤立成分の五直線}
\label{sec:split-point-pencils}
\begin{theorem}\label{thm:split-point-pencils}
可逆二平面分解を持つ全窓は、全位置で左右の射影ファイバーが空または一点である。
一辺の法線・接線の非退化関係を二行列 $C,D$、混合四関係の $NT,TN$ ブロックをそれぞれ
$R_a,S_a$ ($1\le a\le4$) とする。$J=\left(\begin{smallmatrix}0&-1\\1&0\end{smallmatrix}\right)$ とおく。
$n\in\mathbb P(N^*)$ のまわりで座標を一つ1にすれば、$n$ 側の点グラフは
次の四行四列鉛筆と一つの二次式でスキーム論的に与えられる。
\[
 M(n)=\left[\,n^tR_a\ \ (S_a JC^t n)^t\,\right]_{a=1}^4,\qquad
 M(n)\binom{t'}t=0,\qquad t^t D t'=0 .
\]
$t$ は始点の接線座標、$t'$ は終点の接線座標である。
接線側の直線の近傍は $C,D$ と $R_a,S_a$ を交換した式で与えられ、両者で全グラフを覆う。
\end{theorem}
\begin{proof}
始点を $[n:t]$、終点を $[n':t']$ と書く。$n\ne0$ のとき $n'=0$ なら四混合式は
$n\otimes t'$ を零化し、$R_a$ が全二行二列空間を張るので $t'=0$ となる。
従って $n'\ne0$ で、同種式 $n^t Cn'=0$ から、局所座標では $n'=JC^tn$ と規格化できる。
これを他の式へ代入すれば表示の式を得る。鉛筆の左二列は各 $n$ で階数2なので
$t$ を指定すれば $t'$ は高々一つである。交換した近傍でも同じであり、
語順を逆にした混合射影も同型なので反対側も同じである。
これらの消去には単元な小行列しか使わないため、非被約基底でも有効である。
\end{proof}

\begin{proposition}\label{prop:split-simple-four-roots}
両側の座標直線の鉛筆の行列式がそれぞれ相異なる四根を持つとする。
その根で $M(n)$ の核を張る $(t',t)$ が $t^tDt'=0$ を満たすものの数を $g$ とする。
二側でこの数は等しい。被約グラフの像は二本のねじれ直線と、両者を結ぶ
互いに交わらない $g$ 本の直線である。各結線は異なる根どうしを結び、
両端は通常二重点である。残る $4-g$ 個ずつの根では、各々の直線上に長さ1の
埋込み部分があり、これが全ての非被約構造である。
\end{proposition}
\begin{proof}
単根の局所環で鉛筆の階数は3、三つの単元ピボットを除くと最後の式は $zu=0$
（$z$ は根の局所座標）、残りは $h(z)u^2=0$ である。$h(0)$ はまさに述べた核での値である。
$h(0)\ne0$ なら $u^2=zu=0$ で一次元の埋込みイデアル $ku$ を持つ。
$h(0)=0$ なら第二式は第一式から従い、$z=0$ 上の $u$ は任意である。
両方の二列ブロックが単射なので核で $t,t'$ はともに非零である。
従って始点では $[n:ut]$ が座標二直線を結ぶ射影直線の開部分である。
反対の座標直線の近傍でその無限遠点も同じ単根の形となる。
単根の核は一本だけなので、異なる結線は各側で異なる点に至り、
二つの相補平面の直和表示から結線どうしも交わらない。
行列式が単元の点では $t=t'=0$ となるから他の成分はない。
\end{proof}

\begin{corollary}\label{cor:isolated-split-point-scheme}
六次式の孤立成分の唯一の標準窓の点グラフは、二本のねじれ直線と三本の結線を持つ。
六交点は通常二重点で、これとは別にねじれ二直線の上に長さ1の埋込み部分が一つずつある。
特に被約曲線は次数5、算術種数2で、分解の二本の直線は三つの交点を持つ成分として無印で区別できる。
\end{corollary}
\begin{proof}
定理\ref{thm:isolated-split-component}の標数零の代表を使う。
\path{computations/verify_isolated_split_point_scheme.py} は
$\mathbb Q(\sqrt{-1})[d]/((d^2+1)^3-4d^3)$ 上で二切断の相補式から四テンソルを復元する。
各座標直線で $n=(1,y)$ として得る行列式は次数4で、その微分と互いに素である。
随伴行列の四列それぞれを $(t',t)$ として求めた $t^tDt'$ の四式と
行列式との最大公約因子は、両座標直線で次数3である。単根では随伴行列の階数が1のため、
この検査は核による判定と一致する。無限遠で行列式が単元なので全ての点を数えている。
前命題に $g=3$ を適用する。五本の射影直線を六個の通常二重点で貼ると
算術種数は $6-5+1=2$ である。既に最小周期1であるから全位置で同じ記述になる。
\end{proof}

\begin{corollary}\label{cor:no-branch-on-invariant-line-lifts}
標準正則窓に四移動不変の直線グラフがあり、その上の左右評価が全て階数3なら、
この窓のどの辺のどの点にも分岐はない。さらに全点グラフの次元は、
二平面分解への重み特殊化で得られる点グラフの次元以下である。
\end{corollary}
\begin{proof}
定理\ref{thm:invariant-line-filtered-lifts}により重み特殊化の特殊ファイバーは可逆二平面分解になる。
評価階数2以下の点の入射は、階数6の関係束と各位置の射影三空間との積上で閉である。
射影空間は固有なので、族の基底への像も閉である。特殊ファイバーには前定理によりこの点がない。
非零ファイバーは互いに同型なので、どれにも分岐はない。点グラフの次元も射影族の
上半連続性で特殊ファイバーの次元以下となる。
\end{proof}
